\section{Despliegue y monitoreo}

\subsection{Infraestructura de Log + Replay}

Antes de desplegar una offline policy es necesario preparar el logging pipeline:

\begin{itemize}
    \item Sincronizar la salida de la policy con el reward del entorno, para asegurar que el reward trace esté vinculado con el comportamiento de salida.
    \item Registrar la action entropy y la Q variance de la policy, para poder reproducir las fallas.
    \item Definir un alert threshold para los high-risk state (por ejemplo, Q variance > 0.5).
\end{itemize}

\begin{warningbox}{Deployment: el error más común}
Sustituir la robust evaluation por «Offline policy + online fine-tuning»: el reward drift solo se descubre después de salir a producción. Lo correcto es validar primero con un replay simulator y después ir abriendo gradualmente el umbral de online interaction.
\end{warningbox}

\subsection{Panel de monitoreo}

\begin{table}[H]
\centering
\begin{tabular}{p{0.22\textwidth} p{0.26\textwidth} p{0.26\textwidth}}
\toprule
Componente de Dashboard & Métrica & Condición de activación \\
\midrule
Policy behavior & Action distribution shift & KL > 0.3 \\
Reward & Mean reward vs reference & diferencia > 5\% \\
Q stability & Q variance & > 0.4 \\
\bottomrule
\end{tabular}
\caption{Métricas clave a monitorear después del despliegue}
\end{table}

\subsection{Resumen de la sección}

Un despliegue confiable depende del registro de logs, la réplica mediante replay y el panel de monitoreo. Configurar de antemano el alert y el control de entropy permite que la offline policy sea más estable después de salir a producción.

\subsection{Ejemplo de visualización del panel de monitoreo}

\begin{figure}[H]
\centering
\includegraphics[page=4,width=0.9\textwidth]{07_cs224r_offline_rl_2025.pdf}
\caption{El dashboard de monitoreo y los alert thresholds mostrados en las Lecture slides\protect\footnotemark}
\end{figure}
\footnotetext{Página 4 de las Slides, muestra un dashboard típico; la columna central muestra el control de reward drift.}

